\section{Méthode et protocole expérimental}
\subsection{Critère de segmentation}
Pour $x_i\in\R^p$, le coût d'un segment $[a,b]$ est sa somme des carrés
autour du barycentre, suivant le critère d'homogénéité de Fisher
\cite{fisher} :
\begin{equation}
 \mu_{a:b}=\frac{1}{b-a+1}\sum_{i=a}^{b}x_i,
 \qquad D(a,b)=\sum_{i=a}^{b}\norm{x_i-\mu_{a:b}}^2.
 \label{eq:segment}
\end{equation}
Avec $s_0=1<s_1<\dots<s_{K-1}<s_K=n+1$, la partition est formée des
intervalles $[s_{k-1},s_k-1]$ et minimise
\begin{equation}
 C(s_1,\dots,s_{K-1})=\sum_{k=1}^{K}D(s_{k-1},s_k-1).
 \label{eq:objectif}
\end{equation}
Les segments sont non vides, disjoints et couvrent toutes les observations.
Cette écriture conserve l'inertie du clustering euclidien, mais restreint
l'ensemble des partitions admissibles à celles respectant l'ordre.
